\documentclass[11pt, a4paper]{common/document_template}
\usepackage{common}

\begin{document}

\begin{titlepage}
  \begin{center}
	\Huge{Mосковская математическая олимпиада}
  \end{center}
\end{titlepage}

\chapter{2019. Номер 82}

\begin{exercise}
Король вызвал двух мудрецов и объявил им задание: первый задумывает 7 различных натуральных чисел с суммой 100, тайно сообщает их королю, а второму мудрецу называет лишь четвёртое по величине из этих чисел, после чего второй должен отгадать задуманные числа. У мудрецов нет возможности сговориться. Могут ли мудрецы гарантированно справиться с заданием?
\end{exercise}

\begin{Solution}
Заметим, что если четвертое по величине число равно $n$, то сумма чисел больших $n$ будет $\geqslant 3n + 6$, а сумма чисел меньших $n$ будет $\geqslant 6$. То есть сумма всех чисел $100 \geqslant 4n + 12$. Значит если $n = 22$, то все числа определяются однозначно и равны $1, 2, 3, 22, 23, 24, 25$.
\end{Solution}

\end{document}